\begin{abstract}
Updating a language model’s knowledge via fine-tuning or editing is necessary for keeping models up-to-date, yet it often triggers undesired ripple effects such as unintentional forgetting and newly induced hallucinations. Despite growing attention to side effects, we still lack a systematic understanding of which knowledge is most likely to be corrupted and how far the disturbance propagates. We address this gap by constructing a large knowledge topology from a real-world corpus: a verified functional knowledge graph with paired factual/counterfactual triplets, stratified by node in-degree as a proxy for knowledge ``popularity''. We then perform controlled single-fact knowledge updates using parameter-efficient tuning and quantify ripple effects over increasing graph distance using an error propagation rate that captures clean-to-wrong flips.
% and (ii) confidence shift that reveals silent degradation.
Our results show that ripple effects routinely extend well beyond the edited fact, overturning unrelated knowledge even at distances exceeding five hops. Strikingly, high-connectivity hubs, which are widely connected, “popular” knowledge are significantly more likely to flip from correct to incorrect than long-tail knowledge, contradicting the common intuition that popular knowledge is inherently robust. Meanwhile, updates can also distort model calibration: many facts that were wrong before remain wrong after updating but become more confident, producing high-confidence hallucinations. Finally, leveraging graph structure yields an effective mitigation: anchoring updates with a small set of high-connectivity hub facts (selected via our popularity proxy) substantially dampens long-range propagation and reduces hallucinations, outperforming random anchoring, especially at larger distances.


\end{abstract}